% Generated draft: requires mathematical and source review.
\documentclass[11pt]{article}
\usepackage[margin=25mm]{geometry}
\usepackage{fontspec}
\setmainfont{Latin Modern Roman}
\usepackage{amsmath,amssymb,mathtools,unicode-math}
\setmathfont{Latin Modern Math}
\usepackage{enumitem,needspace,xurl,hyperref}
\setlength{\parindent}{0pt}
\setlength{\parskip}{5pt}
\setlength{\emergencystretch}{3em}
\newcommand{\N}{\mathbb{N}}
\newcommand{\Z}{\mathbb{Z}}
\newcommand{\Q}{\mathbb{Q}}
\newcommand{\R}{\mathbb{R}}
\newcommand{\C}{\mathbb{C}}
\newcommand{\F}{\mathbb{F}}
\newcommand{\A}{\mathbb{A}}
\newcommand{\PP}{\mathbb{P}}
\newcommand{\E}{\mathbb{E}}
\newcommand{\eps}{\varepsilon}
\newcommand{\dd}{\,\mathrm{d}}
\begin{document}
\section*{UnsolvedMath ID: 701554}
\textbf{Source problem code:} catalogue entry
\subsection{Definitions and mathematical statement}
Let \(A,B,H\) be groups with specified injective homomorphisms \(i_A:H\hookrightarrow A\) and \(i_B:H\hookrightarrow B\). Their free product with amalgamation is
\[
G=A\ast_H B=(A\ast B)/\langle\!\langle i_A(h)i_B(h)^{-1}:h\in H\rangle\!\rangle.
\]
Identify \(A\), \(B\), and the common copy of \(H\) with their canonical images in \(G\).

A group \(K\) is nilpotent if its lower central series, defined by \(\gamma_1(K)=K\) and \(\gamma_{j+1}(K)=[\gamma_j(K),K]\), reaches the trivial subgroup after finitely many steps. Here \([U,V]\) is the subgroup generated by the commutators \(u^{-1}v^{-1}uv\), with \(u\in U\) and \(v\in V\). A group \(G\) is residually finite if, for every \(g\in G\setminus\{1\}\), there exist a finite group \(Q\) and a homomorphism \(\varphi:G\to Q\) such that \(\varphi(g)\ne1\). The condition that \(H\) satisfies a nontrivial identical relation means that there exist an integer \(n\ge1\) and a nonidentity word \(w\) in the free group on \(n\) generators such that
\[
w(h_1,\ldots,h_n)=1\qquad\text{for every }(h_1,\ldots,h_n)\in H^n.
\]

The supplied record attributes the following two conjectures to M. K. Azarian:
\begin{enumerate}
\item[(a)] For every such amalgam, if both \(A\) and \(B\) are nilpotent, then
\[
\mu(G)\le H.
\]
\item[(b)] For every such amalgam, if \(G\) is residually finite and \(H\) satisfies a nontrivial identical relation, then
\[
\lambda(G)\le H.
\]
\end{enumerate}
\textbf{Unresolved qualification.} The supplied complete record does not define \(\mu(G)\) or \(\lambda(G)\). The notation \(\le\) suggests subgroup containment, but the intended subgroup constructions are unspecified. Consequently, these displays preserve the source assertions but do not constitute fully specified mathematical conjectures until both symbols are defined from reliable source context. The record labels the problem open; this is a supplied status, not an independently verified assessment.
\subsection{Short English statement}
Azarian's two conjectures concern an amalgamated free product: nilpotence of both factors is asserted to imply \(\mu(G)\le H\), while residual finiteness of \(G\) together with a nontrivial group identity on \(H\) is asserted to imply \(\lambda(G)\le H\). The source does not define \(\mu\) or \(\lambda\), so both formulations require clarification.
\subsection{Sources}
\begin{itemize}
\item Local UnsolvedMath registry, permanent ID \(701554\), exact source problem code \texttt{catalogue entry}, title \emph{Kourovka Notebook Problem 19.6}. The supplied record attributes the problem to M. K. Azarian, proposed in 2018, and labels it open.
\item Supplied bibliographic attribution: \emph{The Kourovka Notebook (Unsolved Problems in Group Theory)}, original problem 19.6, 19th issue (2018); the background also identifies the 21st issue edition (2026), \texttt{arXiv:1401.0300v48}. The supplied collection description names E. I. Khukhro and V. D. Mazurov as editors. This attribution and edition information have not been independently verified.
\end{itemize}
\end{document}
